\chapter*{ПРИЛОЖЕНИЕ В. ПРОТОКОЛЫ ЭКСПЕРИМЕНТОВ}
\phantomsection
\addcontentsline{toc}{chapter}{ПРИЛОЖЕНИЕ В. ПРОТОКОЛЫ ЭКСПЕРИМЕНТОВ}

\section*{В.1 Карточка приемочного прогона}
\phantomsection
\addcontentsline{toc}{section}{В.1 Карточка приемочного прогона}

\begin{longtable}{@{}p{0.35\textwidth}p{0.59\textwidth}@{}}
ID прогона & RUN-2026-05-13-A1 \\
Дата и время & 2026-05-13 07:49 UTC \\
Профиль стенда & приемка на CPU-профиле, Docker Compose \\
Команда повторения & \texttt{make e2e-acceptance-artifact} \\
Версия кода & рабочая копия проекта, артефакт \texttt{acceptance\_results\_2026-05-13.json} \\
Конфигурация камер & синтетическая камера, \texttt{1280x720}, \texttt{fps\_target=10} \\
Набор активных правил & \texttt{zone\_enter}; для \texttt{line\_cross} и \texttt{dwell\_time} использованы синтетически зафиксированные события интеграционного контура \\
Длительность прогона & \(30.732\) с \\
Сводный результат & ПРОЙДЕН \\
\end{longtable}

Критерии успешности RUN-2026-05-13-A1 фиксировались до анализа результатов: поле \texttt{result} должно иметь значение \texttt{passed}; функциональное событие должно быть опубликовано и доставлено; при отказе RabbitMQ запись исходящего сообщения должна сохраниться и перейти в состояние \texttt{published} после восстановления брокера; при недоступности HTTP-получателя должна быть зафиксирована повторная доставка; операция записи для роли \texttt{viewer} должна завершиться кодом 403. Промышленная p95-оценка полного видеоконтура в этот критерий не включалась и вынесена в отдельный будущий стендовый прогон.

Артефакты приложения интерпретируются раздельно. Измерительные артефакты содержат значения, которые прямо переносятся в таблицы главы~4: это JSON-файл приемочного прогона и файл \texttt{results.csv} обучения детектора. Проверочные артефакты подтверждают отдельные ветви логики, например переход сообщения в очередь недоставленных сообщений при исчерпании повторов. Демонстрационные артефакты, включая видео БПЛА и снимки интерфейса, подтверждают связность сценария и качество представления результата, но не используются для расчета p95-задержки полного видеоконтура.

\section*{В.2 Протокол метрик прогона}
\phantomsection
\addcontentsline{toc}{section}{В.2 Протокол метрик прогона}

\begin{longtable}{@{}p{0.34\textwidth}p{0.20\textwidth}p{0.38\textwidth}@{}}
\caption{Протокол метрик RUN-2026-05-13-A1}\\
\toprule
Метрика & Значение & Источник данных \\
\midrule
\endfirsthead
\caption*{Продолжение таблицы В.1}\\
\toprule
Метрика & Значение & Источник данных \\
\midrule
\endhead
Задержка функциональной доставки & \(4.605\) с & поле \texttt{functional\_delivery} блока \texttt{timings\_sec} \\
Задержка восстановления RabbitMQ & \(8.863\) с & поле \texttt{rabbitmq\_recovery} блока \texttt{timings\_sec} \\
Задержка восстановления HTTP-получателя & \(4.455\) с & поле \texttt{webhook\_recovery} блока \texttt{timings\_sec} \\
Повтор при отказе RabbitMQ & 1 & поле \texttt{retry\_count} в блоке \texttt{rabbitmq\_reliability} \\
Успешные доставки, функциональный сценарий & 1 & счетчик \texttt{success} в блоке \texttt{functional\_e2e} \\
Неуспешные доставки, функциональный сценарий & 0 & счетчик \texttt{failed} в блоке \texttt{functional\_e2e} \\
Неуспешные доставки, повтор HTTP-доставки & 2 & счетчик \texttt{failed} в блоке \texttt{webhook\_retry} \\
Число сообщений в очереди недоставленных & 0 & приемочный артефакт; штатный прогон без конечной ошибки доставки \\
Статус создания камеры ролью viewer & 403 & проверка ограничения доступа \\
\bottomrule
\end{longtable}

\section*{В.3 Протокол отказных сценариев}
\phantomsection
\addcontentsline{toc}{section}{В.3 Протокол отказных сценариев}

\begin{longtable}{@{}p{0.35\textwidth}p{0.59\textwidth}@{}}
Сценарий отказа & Кратковременное отключение RabbitMQ при наличии события в таблице исходящих сообщений \\
Потеря событий & 0 \\
Статус во время отказа & \texttt{retry}, \texttt{retry\_count=1}, \texttt{published\_at} пустой \\
Статус после восстановления & \texttt{published}, \texttt{retry\_count=1}, \texttt{published\_at=2026-05-13T07:49:52+00} \\
Время восстановления & \(8.863\) с \\
Состояние очереди недоставленных сообщений & сообщений не зафиксировано \\
Комментарий & Таблица исходящих сообщений сохранила событие до восстановления брокера; доставка завершилась без повторной генерации события аналитическим контуром \\
\end{longtable}

\begin{longtable}{@{}p{0.35\textwidth}p{0.59\textwidth}@{}}
Сценарий отказа & Временная недоступность HTTP-получателя \\
Потеря событий & 0 \\
Неуспешные попытки до восстановления & 1 зафиксированная попытка до успешной доставки; всего по событию: success \(=1\), failed \(=2\) \\
Итоговый статус исходящего сообщения & \texttt{published}, \texttt{published\_at=2026-05-13T07:49:54+00} \\
Время восстановления & \(4.455\) с \\
Состояние очереди недоставленных сообщений & сообщений не зафиксировано \\
Комментарий & Политика повторной доставки позволила доставить событие после восстановления \texttt{webhook-mock} \\
\end{longtable}

\section*{В.4 Сохраненные артефакты}
\phantomsection
\addcontentsline{toc}{section}{В.4 Сохраненные артефакты}

\begin{itemize}
    \item JSON-артефакт приемочного прогона RUN-2026-05-13-A1: \path{docs/MasterWork/REPORT/LaTeX/acceptance_results_2026-05-13.json};
    \item таблица метрик обучения: \path{runs/detect/runs/train/small_detection/results.csv};
    \item график обучения: \path{runs/detect/runs/train/small_detection/results.png}, включен в отчет как \path{figures/yolo_training_results.png};
    \item матрица ошибок: \path{runs/detect/runs/train/small_detection/confusion_matrix.png}, включена в отчет как \path{figures/yolo_confusion_matrix.png};
    \item итоговые веса детектора: \path{runs/detect/runs/train/small_detection/weights/best.pt};
    \item демонстрационное видео: \path{videos/V_DRONE_009.mp4};
    \item тест логики повторной доставки и очереди недоставленных сообщений: \path{services/integration-worker/tests/test_consumer_retry.py};
    \item скриншоты пользовательского интерфейса: \path{figures/ui-dashboard.png}, \path{figures/ui-demo.png}, \path{figures/ui-status.png}, \path{figures/ui-events.png}.
\end{itemize}

\section*{В.5 Чек-лист функциональных сценариев}
\phantomsection
\addcontentsline{toc}{section}{В.5 Чек-лист функциональных сценариев}

\scriptsize
\begin{longtable}{@{}p{0.07\textwidth}p{0.27\textwidth}p{0.39\textwidth}p{0.19\textwidth}@{}}
\caption{Чек-лист функционального тестирования}\\
\toprule
ID & Сценарий & Ожидаемый результат & Статус \\
\midrule
\endfirsthead
\caption*{Продолжение таблицы В.2}\\
\toprule
ID & Сценарий & Ожидаемый результат & Статус \\
\midrule
\endhead
F-01 & Вход c валидными учетными данными & Выдача маркеров доступа и обновления & ПРОЙДЕН \\
F-02 & Обновление токена & Ротация маркера доступа без повторного входа & ПРОЙДЕН \\
F-03 & Завершение сессии & Маркер обновления отозван, повторное обновление невозможно & ПРОЙДЕН \\
F-04 & Создание камеры & Запись добавлена, камера доступна в списке & ПРОЙДЕН \\
F-05 & Обновление параметров камеры & Изменения сохраняются, доступны через API и подхватываются фоновым обработчиком & ПРОЙДЕН \\
F-06 & Создание зоны polygon & Геометрия валидирована и сохранена & ПРОЙДЕН \\
F-07 & Создание правила zone\_enter & Правило активно и доступно фоновому обработчику & ПРОЙДЕН \\
F-08 & Сценарий zone\_enter & Событие создано один раз и доставлено & ПРОЙДЕН \\
F-09 & Сценарий line\_cross & Событие проходит через таблицу исходящих сообщений и HTTP-получатель & ПРОЙДЕН, синтет. \\
F-10 & Сценарий dwell\_time & Событие проходит через контур повторной HTTP-доставки & ПРОЙДЕН, синтет. \\
F-11 & Дедупликация по dedup\_key & Повторное сообщение о событии не создает дубликат факта & ПРОЙДЕН, схема \\
F-12 & Просмотр журнала событий & Фильтрация по камере, типу и времени поддержана пользовательским интерфейсом и API & ПРОЙДЕН \\
\bottomrule
\end{longtable}
\normalsize

\section*{В.6 Чек-лист надежности и безопасности}
\phantomsection
\addcontentsline{toc}{section}{В.6 Чек-лист надежности и безопасности}

\scriptsize
\begin{longtable}{@{}p{0.07\textwidth}p{0.27\textwidth}p{0.39\textwidth}p{0.19\textwidth}@{}}
\caption{Чек-лист надежности и безопасности}\\
\toprule
ID & Сценарий & Ожидаемый результат & Статус \\
\midrule
\endfirsthead
\caption*{Продолжение таблицы В.3}\\
\toprule
ID & Сценарий & Ожидаемый результат & Статус \\
\midrule
\endhead
R-01 & Отключение RabbitMQ при генерации событий & События остаются в таблице исходящих сообщений, потерь нет & ПРОЙДЕН \\
R-02 & Восстановление RabbitMQ & Накопленная очередь исходящих сообщений обрабатывается, публикация возобновляется & ПРОЙДЕН \\
R-03 & Ошибка HTTP-получателя 5xx/недоступность & Срабатывает политика повторной доставки & ПРОЙДЕН \\
R-04 & Исчерпание лимита повторов & Сообщение попадает в очередь недоставленных & ЧАСТИЧНО, тест логики \\
R-05 & Недоступность одной камеры & Другие камеры продолжают обработку & СЛЕДУЮЩИЙ СТЕНД \\
S-01 & Роль viewer пытается изменить камеру & Доступ запрещен (403) & ПРОЙДЕН \\
S-02 & Роль operator создает/меняет правила & Операция разрешена & ПРОЙДЕН \\
S-03 & Роль admin управляет пользователями & Операция разрешена & ПРОЙДЕН \\
S-04 & Запрос без JWT & Доступ к защищенному API запрещен & ПРОЙДЕН \\
\bottomrule
\end{longtable}
\normalsize

\section*{В.7 Итоговый акт стендового прогона}
\phantomsection
\addcontentsline{toc}{section}{В.7 Итоговый акт стендового прогона}

\begin{longtable}{@{}p{0.35\textwidth}p{0.59\textwidth}@{}}
Номер акта & ACT-RUN-2026-05-13-A1 \\
Дата и время & 2026-05-13 07:49 UTC \\
Конфигурация стенда & приемка на CPU-профиле, Docker Compose, PostgreSQL, RabbitMQ, API, analytics-worker, integration-worker, webhook-mock \\
Номер версии & рабочая копия проекта, приемочный артефакт сохранен в каталоге отчета \\
Набор сценариев & Интеграционная доставка, отказ и восстановление RabbitMQ, повторная HTTP-доставка, RBAC-ограничение для роли \texttt{viewer} \\
Сводный результат & ПРОЙДЕН \\
Зафиксированные отклонения & промышленная p95-оценка для многокамерного профиля не выполнялась; требуется отдельная длительная серия прогонов \\
Решение по приемке & Платформа принята как воспроизводимый исследовательско-экспериментальный стенд текущего этапа \\
Ответственный исполнитель & Нурмухамедов А.\,Б. \\
\end{longtable}
